\begin{figure}[H]
\centering
\begin{subfigure}[t]{0.47\textwidth}
\centering
\begin{tikzpicture}
\begin{axis}[width=\linewidth,height=5.8cm,xlabel={$x$},
  xmin=-4,xmax=6,ymin=0,ymax=1.3,
  legend style={font=\scriptsize,at={(0.5,-0.28)},anchor=north,legend columns=2},ytick=\empty,axis lines=left]
  \addplot[uniblue,very thick,domain=-4:6,samples=200] {exp(-x^2/2)};
  \addlegendentry{$\psi(x)$}
  \addplot[unired,very thick,domain=-4:6,samples=200] {exp(-(x-2)^2/2)};
  \addlegendentry{$\hat T(a)\psi=\psi(x-a)$}
  \draw[-{Stealth},thick] (axis cs:0,1.08) -- (axis cs:2,1.08)
        node[midway,above] {$a>0$};
\end{axis}
\end{tikzpicture}
\caption{Con $a=2$: la funci\'on se desplaza hacia la \textbf{derecha}.}
\end{subfigure}\hfill
\begin{subfigure}[t]{0.52\textwidth}
\centering
\begin{tikzpicture}
\begin{axis}[width=\linewidth,height=5.8cm,xlabel={$x$},
  xmin=-3,xmax=4,ymin=-0.9,ymax=1.5,
  legend style={font=\scriptsize,at={(0.5,-0.28)},anchor=north,legend columns=3},
  grid=both,minor tick num=1]
  \addplot[black,very thick,domain=-3:4,samples=200] {exp(-(x-1)^2/2)};
  \addlegendentry{exacta $\psi(x-1)$}
  \addplot[uniblue,dashed,thick,domain=-3:4,samples=200] {exp(-x^2/2)};
  \addlegendentry{$N=0$}
  \addplot[unired,thick,domain=-3:4,samples=200] {exp(-x^2/2)*(1+x)};
  \addlegendentry{$N=1$}
  \addplot[unigreen,thick,domain=-3:4,samples=200]
    {exp(-x^2/2)*(1+x+(x^2-1)/2)};
  \addlegendentry{$N=2$}
  \addplot[orange,thick,domain=-3:4,samples=200]
    {exp(-x^2/2)*(1+x+(x^2-1)/2+(x^3-3*x)/6)};
  \addlegendentry{$N=3$}
\end{axis}
\end{tikzpicture}
\caption{Sumas parciales de $\sum_{n}\frac{(-a)^n}{n!}\psi^{(n)}(x)$, con $a=1$ y $\psi=\eu^{-x^2/2}$.}
\end{subfigure}
\caption{Acci\'on de $\hat T(a)=\eu^{-a\,\dif/\dif x}$: el desarrollo de Taylor reconstruye $\psi(x-a)$.}
\label{fig:traslacion}
\end{figure}